The discovery of REM sleep enabled the first major \textit{neurobiological} theory of dreaming, which was formulated as a direct rebuttal to Freudian psychology. In 1977, \citet{hobson1977activation} proposed the "Activation-Synthesis Hypothesis" (ASH).

ASH proposed that the "bizarreness" of dreams was not symbolic (contra Freud) but mechanical. The mechanism was twofold:
\begin{enumerate}
    \item \textbf{Activation:} During REM sleep, the brainstem (specifically the pons) becomes active, sending \textit{random} ponto-geniculo-occipital (PGO) waves up to the forebrain.
    \item \textbf{Synthesis:} The higher-level forebrain, or cortex, is then activated by these random signals. The cortex, which is "always trying to make sense of" input, "synthesizes" these random signals into a coherent narrative, "making the best of a bad job".
\end{enumerate}

In this model, the dream is simply the forebrain's "attempt to make sense of random signals", and its bizarreness is due to a "loss of the organizing capacity of the brain," not an "elaborate disguise mechanism" \citep{hobson1977activation}.

The Activation-Synthesis Hypothesis was a critical step, shifting the focus from psychology to neurobiology. While its specific claims of "randomness" are now considered outdated, ASH serves as the direct intellectual ancestor to the modern computational models described in Section 3.

Hobson and McCarley correctly identified the \textit{mechanism}: a two-part process involving a bottom-up signal generator (Activation) and a top-down interpreter that builds a world. The computational theories of today are a \textit{refinement} of this exact idea, not a replacement.